\documentclass[10pt,a4paper]{amsart}
\usepackage[margin=19mm]{geometry}
\usepackage{amsmath,amssymb,mathtools}
\usepackage[scr=rsfs,cal=euler]{mathalfa}
\usepackage{xfrac}
\usepackage[unicode,hidelinks,bookmarks=false]{hyperref}
\newcommand{\ie}{i.e.\ }
\newcommand{\cf}{cf.\ }
\newcommand{\resp}{respectively\ }
\newcommand{\Subsection}{\subsection}
\providecommand{\nobreakdash}{\nobreak\hbox{-}}
\setlength{\emergencystretch}{2em}
\multlinegap0pt
\usepackage{amssymb}
\usepackage{mathtools}
\usepackage{bm}
\usepackage{xcolor}
\usepackage{csquotes}
\usepackage{tikz}
\newtheorem{lemma}{Lemma}
\newcommand{\glang}{\Lambda}
\newcommand{\univ}{\Sigma}
\title{The Polymatroid Representation of a Greedoid, and Associated Galois Connections}
\author{Robert P. Streit, Vijay K. Garg}
\date{Source: arXiv:2411.15363v4 (CC BY 4.0)}
\begin{document}
\maketitle

\noindent\emph{Source and licence.} The Polymatroid Representation of a Greedoid, and Associated Galois Connections. Version arXiv:2411.15363v4 (CC BY 4.0), distributed by arXiv under the Creative Commons Attribution 4.0 International licence, http://creativecommons.org/licenses/by/4.0/, as recorded in the arXiv licence metadata for this exact version and shown on the abstract page https://arxiv.org/abs/2411.15363v4. Full licence notice: \texttt{licenses/arxiv-2411.15363-CC-BY-4.0.txt}.\par\smallskip

\noindent\textit{Editorial scope.} Counted unit: the article's lemma lem:forking (source0.tex lines 1383--1387). The article's complete original proof of it is delivered with it (source0.tex lines 1428--1480, one \texttt{proof} environment of 53 lines, which the article places away from the statement). No mathematics is written, repaired or re-derived: the statement and the proof are the article's own lines, in the article's own notation, and no retained line is substituted or re-flowed.  Retained uncounted as the source-local context the proof needs: lines 1398--1401.  Nothing was omitted from any retained span, so the package carries no omission record.  No retained line contains a citation, so the wrapper carries no bibliography and no bibliography entry.  No retained line contains a figure environment, an image-inclusion command or drawing code, so no image is delivered and the package declares no asset dependency.  Editorial formatting only, in addition: an AMS \texttt{amsart} preamble that restates the article's own declarations and macros used by the retained lines, namely the environments \texttt{lemma} on separate counters, together with the standard packages those lines need.  The retained environments print in this document as Lemma 1 in order of appearance, whereas the article prints the same items with the numbering of its own full document.  The complete native source of the article as distributed in the arXiv e-print for this exact version is bundled unchanged as \texttt{sources/168980/source0.tex}.
\begin{lemma}\label{lem:inner-forking}
    Let $\glang$ be a greedoid with the interval property, and $\mu \in \glang$ a feasible word.
    If there exists $x,y \in \Gamma[\mu]$ and $\alpha \in \univ^*$ with $\mu x \alpha$ and $\mu y \alpha$ feasible, then $\mu x \alpha \not\sim \mu y \alpha$ implies $\mu y \alpha x \in \glang$.
\end{lemma}

\begin{lemma}[Forking Lemma]\label{lem:forking}
    Let $\glang$ be a greedoid with the interval property, and select two flats $F, F'\in\glang/\mathord{\sim}$.
    Suppose that $\mu \in F\sqcap F'$.
    Then, $x \in \Gamma(F\sqcap F')$ and $[\mu x] \sqsubseteq F'$ implies $x \in \Gamma(F)$.
\end{lemma}

\begin{proof}[Proof of the Forking Lemma]
    Let $\mu \in \glang$ and $x \in \Gamma[\mu]$ be as described in Lemma~\ref{lem:forking}.
    Put more concisely, our assumptions make,
    $$F\sqcap F' = [\mu] \sqsubset [\mu x] \sqsubseteq F'.$$
    and from this, the reader should observe that the strict ordering relation between $[\mu]$ and $[\mu x]$ ensures $[\mu x] \not\sqsubseteq F$, thereby making $F \neq F'$ as well.
    Hence, there exists a letter $y \in \Gamma[\mu]$ and $\alpha \in \univ^*$ where $y \neq x$ and $\mu y \alpha \in F$.
    To prove the lemma, we will induct over the difference $r(F) - r[\mu]$.
    In the base case where $r(F) - r(F\sqcap F') = 1$, we see that $\alpha = \epsilon$.
    Then, because $\mu y \in F$ and $[\mu x] \not \sqsubseteq F$, it is the case that $\mu x \not\sim \mu y$.
    Therefore, by applying Lemma~\ref{lem:inner-forking} to $\mu x \epsilon$ and $\mu y \epsilon$ we find $\mu y \epsilon x = \mu y x$ is feasible, which makes the base case as this is equivalent to stating $x \in \Gamma(F)$.

    Now, recalling $\mu y \alpha \in F$, for the inductive step we make $r(F) - r[\mu] > 1$ by assuming $\alpha = \beta z$ where $z \in \univ$ is a letter and $\beta \in \univ^*$ is a simple word, possibly $\beta = \epsilon$.
    By the interval property there exists a subword of $\mu y \beta z$ of length at least $|\mu y \beta z| - |\mu x| = |\beta| + 1$ which can be concatenated onto $\mu x$ to make a feasible word.
    The resulting feasible word must be simple (no letters can be repeated), and so there are three possibilities for the identity of this observed subword:
    \begin{enumerate}
        \item The subword equals $y\beta$, and so $\mu x y \beta$ is feasible,
        \item The subword equals $\beta z$, and so $\mu x \beta z$ is feasible,
        \item There exists another subword $\beta'$ of $\beta$ resulting from the deletion of some single letter, which is to say $|\beta'| = |\beta| -1$, such that $\mu x y \beta' z$ is feasible.
    \end{enumerate}
    Suppose that either (1.) or (3.) is true.
    Then, in either of these cases the hereditary axiom makes $\mu xy$ feasible, and thus $\mu y x$ is feasible as well by exchange from $\mu xy$ to $\mu y$.
    Furthermore, $[\mu x] \not\sqsubseteq F$ implies $[\mu y x] \not\sqsubseteq F$ as well (since $[\mu x] \prec [\mu y x]$ by $y \in \Gamma[\mu x]$), and so $[\mu y] \sqsubseteq F$ makes $F \sqcap [\mu yx] = [\mu y]$ by construction.
    As a consequence,
    \[
        r(F) - r(F \sqcap [\mu y x]) = r(F) - r([\mu y]) = r(F) - r([\mu]) + 1.
    \]
    Therefore, by applying the inductive hypothesis to $F$ and $[\mu y x]$, we witness $x \in \Gamma(F)$.
    The only remaining case is (2.), where $\mu x \beta z \in \glang$.
    Here, we simply apply Lemma~\ref{lem:inner-forking} to $\mu x \beta z$ and $\mu y \beta z$ to find $\mu y \beta z x \in \glang$.
    Then, because this case makes $\mu y \beta z \in F$ we conclude that $x \in \Gamma(F)$.
\iffalse
    For the other direction, suppose that for all choices of $F,F' \in \glang/\mathord{\sim}$ we have the following: If a continuation $x \in \Gamma(F\sqcap F')$ satisfies $(F\sqcap F')\circ x \sqsubseteq F'$, then $x \in \Gamma(F)$.
    Select any feasible $\alpha, \beta \in \glang$, and assume $|\beta| > |\alpha|$.
    We will show the existence of a subword $\beta'$ of $\beta$ with length $|\beta'| \geq \delta$, for $\delta \triangleq |\beta| - |\alpha|$, such that $\alpha\beta'$ is feasible.
    To do so, first let $\beta = b_1\ldots b_k$, and let $b_1 \ldots b_i$ be the $\beta$-prefix spanning $[\alpha]\sqcap [\beta]$.
    Note, in the case where $[\alpha] \sqcap [\beta]$ is the bottom flat this prefix equals the empty word $\epsilon$, and we let $i = 0$ by convention.
    We claim that $\beta' = b_{i+1}\ldots b_{i+\delta}$ is a suitable subword for showing the interval property, i.e. that $\alpha b_{i+1}\ldots b_{i+\delta}$ is feasible.
    To show this, we will verify that
    \begin{align*}
        \alpha b_{i+1} \ldots b_{i+j}\in\glang,
    \end{align*}
    for each $j \in \{1, \ldots, \delta\}$.
    We use induction.
    The base case, where $j = 1$, follows easily:
    By construction $b_{i+1} \in \Gamma([\alpha]\sqcap[\beta])$ and,
    \[
        ([\alpha]\sqcap[\beta])\circ b_{i+1} \sqsubseteq [\beta],
    \]
    and so our initial assumptions make $\alpha b_{i+1}$ feasible.
    Now, let $j > 1$ but less than $\delta$.
    By hypothesis, we find that $\alpha b_{i+1}\ldots b_{i + j - 1}$ is feasible.
\fi
\end{proof}

\end{document}
